\section{Introduction}

Generative artificial intelligence has expanded the range of engineering activities that
can be supported by automated information processing. Large language models can interpret
natural-language specifications, generate technical artifacts, support model creation, and
assist with the analysis of heterogeneous engineering information. Recent work has shown
both the potential of such approaches and their limitations in systems engineering and
model-based engineering contexts \cite{Topcu2025,Stein2026,Cibrian.2025}.

The introduction of generative AI into engineering workflows is nevertheless not equivalent
to the establishment of a reliable engineering process. Engineering information is often
distributed across heterogeneous legacy sources, differs in structure, terminology,
granularity, and abstraction level, and may evolve over time
\cite{Maheshwari2018,Henderson2024,Legat2014}. In such environments, technical access to
information does not by itself resolve the question of what the information means, how
reliable a particular interpretation is, or whether it may be used as the basis for a
subsequent engineering decision.

This distinction is particularly relevant for probabilistic AI. A generated result may be
linguistically convincing or formally valid while still containing engineering-relevant
errors or unsupported assumptions \cite{Topcu2025,Stein2026}. Consequently, the central
question is not only whether AI can produce engineering-related output, but under which
conditions such output can become part of a controlled engineering workflow.

The investigation underlying this paper addressed this problem in the context of a
controlled AI-supported transformation of heterogeneous legacy engineering information
into SysML v2. The developed architecture combined source-grounded information,
probabilistic interpretation, explicit human authorization, controlled downstream
transformation, provenance, traceability, and technical validation. Repeated transformation
runs were used to refine the boundaries between information processing, engineering
decision-making, and model formation.

Based on these findings and the literature used in the underlying investigation, this paper
derives generalized design principles for AI-supported engineering. The focus is deliberately
not placed on one particular technical implementation. Instead, the paper addresses the
questions that need to be resolved when probabilistic AI becomes part of an engineering
process: what information may support engineering statements, where interpretation is
required, how uncertainty is handled, where engineering authority resides, how processing
is constrained, and how resulting artifacts remain traceable and subject to assessment
distinct from their generation.

The remainder of the paper is structured as follows. Section II outlines the main challenges
of AI-supported engineering. Section III derives design principles for controlled use.
Section IV organizes these principles into a structured path toward productive engineering
application. Section V discusses transferability, context dependence, and limitations, and
Section VI concludes the paper.
